%% kapa: exam=1 round=35 year=2024 period=2 form=A subject=accounting subject_ko=회계학 seq=1 number=41
%% source: Q-Net 감정평가사 기출문제 (cst00309) artlSeq=5213229 file=raw/1008/5213229/2교시 문제지 원본.pdf
%% conversion: kordoc 4.24.51 -> tools/gen1.py (auto); review=proofread-sonnet; flags=none
\begin{question}{41}
재무보고를 위한 개념체계에 관한 설명으로 옳지 않은 것은?
\begin{choices}
  \item 경제적효익의 유입가능성이나 유출가능성이 낮더라도 자산이나 부채가 존재할 수 있다.
  \item 부채가 발생하거나 인수할 때의 역사적 원가는 발생시키거나 인수하면서 수취한 대가에서 거래원가를 가산한 가치이다.
  \item 매각이나 소비되는 자산의 원가에 대한 정보와 수취한 대가에 대한 정보는 예측가치를 가질 수 있다.
  \item 가격 변동이 유의적일 경우, 현행원가를 기반으로 한 이익은 역사적 원가를 기반으로 한 이익보다 미래 이익을 예측하는데 더 유용할 수 있다.
  \item 합리적인 추정의 사용은 재무정보 작성의 필수적인 부분이며 추정치를 명확하고 정확하게 기술하고 설명한다면 정보의 유용성을 훼손하지 않는다.
\end{choices}
\end{question}
\answer{2}
